\cleardoublepage
\phantomsection
\begin{planbox}[แผนบริหารการสอนประจำบทที่ 7\quad ไฟฟ้าสถิต ประจุผิวเซลล์ และการย้อมสีแกรม (สัปดาห์ที่ 8)]
\noindent\textbf{1. หัวข้อเนื้อหาประจำบท (บรรยาย 2 ชั่วโมง และปฏิบัติการ 3 ชั่วโมง)}
\begin{itemize}
    \item 7.1 ประจุไฟฟ้า กฎของคูลอมบ์ และแรงทางไฟฟ้าระหว่างประจุบวกและประจุลบ
    \item 7.2 สนามไฟฟ้า ศักย์ไฟฟ้า และความต่างศักย์เยื่อหุ้มเซลล์จุลินทรีย์
    \item 7.3 กลไกทางไฟฟ้าสถิตของการย้อมสีแบคทีเรียแบบธรรมดา (Simple Staining) การย้อมสีแกรม (Gram Staining) และการย้อมสีพื้นหลัง (Negative Staining)
    \item 7.4 ปฏิบัติการที่ 7 การศึกษาแรงทางไฟฟ้าสถิตและเปรียบเทียบการจับตัวของสีย้อมประจุบวกกับสีย้อมประจุลบ
\end{itemize}

\noindent\textbf{2. วัตถุประสงค์เชิงพฤติกรรมและผลลัพธ์การเรียนรู้รายวิชา (CLOs ที่เชื่อมโยง PLOs)}
\begin{enumerate}
    \item อธิบายกฎของคูลอมบ์และคำนวณขนาดของแรงทางไฟฟ้าระหว่างประจุได้ \textbf{(CLO1 เชื่อมโยง PLO2)}
    \item อธิบายเหตุผลทางฟิสิกส์ว่าทำไมผนังเซลล์แบคทีเรียที่ค่า pH เป็นกลางจึงมีประจุสุทธิเป็นลบ และทำไมสีย้อมพื้นฐาน เช่น Crystal Violet และ Methylene Blue จึงต้องมีประจุบวก \textbf{(CLO1, CLO2 เชื่อมโยง PLO2, PLO3)}
    \item ปฏิบัติการทดลองเรื่องไฟฟ้าสถิตและเปรียบเทียบการยึดเกาะของสีย้อมอย่างปลอดภัยและเป็นทีม \textbf{(CLO3, CLO4 เชื่อมโยง PLO1, PLO5, PLO6)}
\end{enumerate}

\noindent\textbf{3. กิจกรรมการจัดการเรียนรู้เชิงรุก (Active Learning \& OBE)}
\begin{itemize}
    \item \textbf{Visual Concept Mapping}\quad เชื่อมโยงโครงสร้างกรดไทโคอิก (Teichoic Acid) และลิโพพอลิแซ็กคาไรด์ (LPS) ที่มีหมู่ฟอสเฟต ($\text{PO}_4^{3-}$) บนผิวแบคทีเรีย เข้ากับแรงดึงดูดทางไฟฟ้าสถิตของสีย้อมประจุบวก
\end{itemize}

\noindent\textbf{4. สื่อการเรียนรู้และแหล่งเรียนรู้}
\begin{itemize}
    \item เอกสารคำสอนบทที่ 7 ชุดทดลองไฟฟ้าสถิต และชุดสีย้อม Methylene Blue (ประจุบวก) กับ Nigrosin/Eosin (ประจุลบ)
\end{itemize}

\noindent\textbf{5. การวัดและประเมินผลประจำบท}
\begin{itemize}
    \item แบบทดสอบย่อยท้ายคาบเรื่องกฎของคูลอมบ์และประจุของสีย้อม \textbf{(CLO1, CLO2)}
    \item รายงานปฏิบัติการที่ 7 \textbf{(CLO3, CLO4)}
\end{itemize}
\end{planbox}
\clearpage

\chapter{ไฟฟ้าสถิต ประจุผิวเซลล์ และการย้อมสีแกรม}
\index{ไฟฟ้าสถิต}\index{กฎของคูลอมบ์}\index{ประจุผิวเซลล์แบคทีเรีย}\index{การย้อมสีแกรม}

\section{ประจุไฟฟ้าและกฎของคูลอมบ์}

\deepskinfig{figures/sec_3d/ch07_s1.jpg}{แรงดึงดูดทางไฟฟ้าสถิตตามกฎของคูลอมบ์ ($F_e = k|q_1 q_2|/r^2$) ระหว่างไอออนสีย้อมประจุบวกกับหมู่ฟอสเฟตประจุลบบนผนังเซลล์แบคทีเรีย}

ในธรรมชาติ ประจุไฟฟ้าแบ่งออกเป็น 2 ชนิด ได้แก่ \textbf{ประจุบวก} (\textbf{Positive Charge}, $+q$) และ\textbf{ประจุลบ} (\textbf{Negative Charge}, $-q$) มีหน่วยเอสไอคือ \textbf{คูลอมบ์} ($\text{C}$) กฎพื้นฐานทางไฟฟ้าสถิตระบุว่า \textbf{ประจุชนิดเดียวกันจะผลักกัน และประจุต่างชนิดกันจะดึงดูดกัน} โดยขนาดของแรงทางไฟฟ้า ($F_e$) ระหว่างประจุ $q_1$ และ $q_2$ ที่อยู่ห่างกันเป็นระยะทาง $r$ เป็นไปตาม\textbf{กฎของคูลอมบ์} (\textbf{Coulomb's Law}) ดังแสดงใน\mbox{ภาพที่} 7.1 และสมการที่ (7.1)

\begin{equation}
    F_e = \frac{k |q_1 q_2|}{r^2}
\end{equation}

โดยที่ $k \approx 9.0 \times 10^9\text{ N}\cdot\text{m}^2/\text{C}^2$ คือค่าคงตัวของคูลอมบ์ในสุญญากาศ สังเกตจากสมการได้ว่า แรงทางไฟฟ้าจะแปรผกผันกับกำลังสองของระยะห่าง ($F_e \propto 1/r^2$) กล่าวคือ ถ้าระยะห่างระหว่างโมเลกุลสีย้อมกับผนังเซลล์ลดลงครึ่งหนึ่ง แรงดึงดูดทางไฟฟ้าจะเพิ่มสูงขึ้นถึง $2^2 = 4$ เท่า

\section{ฟิสิกส์ของการย้อมสีแบคทีเรีย ทำไมผนังเซลล์จึงดึงดูดสีย้อมประจุบวก}

\deepskinfig{figures/sec_3d/ch07_s2.jpg}{เปรียบเทียบกลไกการย้อมสีแบคทีเรียด้วยสีย้อมประจุบวก (Basic Dye ดูดติดผนังเซลล์ประจุลบ) และสีย้อมประจุลบ (Acidic Dye ถูกผลักออกเพื่อดูแคปซูลใส)}

ที่สภาวะความเป็นกรด-ด่างเป็นกลาง ($\text{pH } 7.0$) ผนังเซลล์และเยื่อหุ้มเซลล์ของแบคทีเรียเกือบทุกชนิดจะมี\textbf{ประจุไฟฟ้าสุทธิเป็นลบ} (\textbf{Net Negative Surface Charge}) เนื่องจากมีหมู่ฟอสเฟต ($\text{PO}_4^{3-}$) และหมู่คาร์บอกซิล ($\text{COO}^-$) อยู่ในโครงสร้างของกรดไทโคอิก (\textbf{Teichoic Acid} ในแบคทีเรียแกรมบวก) และลิโพพอลิแซ็กคาไรด์ (\textbf{Lipopolysaccharide} ในแบคทีเรียแกรมลบ) รวมถึงสายดีเอ็นเอและอาร์เอ็นเอภายในเซลล์ก็มีหมู่ฟอสเฟตที่มีประจุลบเช่นเดียวกัน

ด้วยหลักการแรงดึงดูดระหว่างประจุต่างชนิดตามกฎของคูลอมบ์ นักจุลชีววิทยาจึงแบ่งเทคนิคการย้อมสีออกเป็น 2 กลไกทางฟิสิกส์ ดังแสดงใน\mbox{ภาพที่} 7.2 และ\mbox{ภาพที่} 7.3

\begin{figure}[H]
\centering
\resizebox{0.88\textwidth}{!}{%
\begin{tikzpicture}[font=\sffamily, >=Stealth]
    % Left: Basic Dye (Positive Charge) -> Direct Staining
    \begin{scope}[shift={(0,0)}]
        \draw[fill=boxBlueBg, draw=physRoyalBlue, line width=1.2pt, rounded corners=4pt] (0,0) rectangle (6.8,4.6);
        \node[font=\bfseries\color{physRoyalBlue}] at (3.4,4.15) {(ก) สีย้อมประจุบวก (Basic Dye)};
        \draw[fill=labPurple!25, draw=labPurple, line width=1.8pt, rounded corners=12pt] (1.6,1.6) rectangle (5.2,3.0);
        \node[font=\bfseries\color{labPurple}] at (3.4,2.3) {เซลล์แบคทีเรีย (ผิวประจุลบ $-$)};
        \foreach \x in {1.9, 2.6, 3.4, 4.2, 4.9} {
            \node[font=\small\bfseries, color=crimsonAlert] at (\x,3.2) {$-$};
            \node[font=\small\bfseries, color=crimsonAlert] at (\x,1.4) {$-$};
        }
        % Positive Dye Ions attracted
        \foreach \x in {2.1, 3.4, 4.7} {
            \fill[physCyan] (\x,3.7) circle (4.5pt);
            \node[white, font=\scriptsize\bfseries] at (\x,3.7) {$+$};
            \draw[->, line width=1.1pt, color=physRoyalBlue] (\x,3.55) -- (\x,3.3);
        }
        \node[font=\footnotesize, align=center, color=charcoal] at (3.4,0.6) {เช่น Crystal Violet, Methylene Blue, Safranin\\ประจุ $+$ ดูดติดกับผิวเซลล์ประจุ $-$ $\rightarrow$ \textbf{ตัวเซลล์ติดสีชัด}};
    \end{scope}

    % Right: Acidic Dye (Negative Charge) -> Negative Staining
    \begin{scope}[shift={(7.6,0)}]
        \draw[fill=boxAmberBg, draw=physAmber, line width=1.2pt, rounded corners=4pt] (0,0) rectangle (6.8,4.6);
        \node[font=\bfseries\color{physAmber!90!black}] at (3.4,4.15) {(ข) สีย้อมประจุลบ (Acidic Dye)};
        \fill[charcoal!75, rounded corners=3pt] (0.9,1.2) rectangle (5.9,3.4);
        \draw[fill=white, draw=physAmber, line width=1.6pt, rounded corners=12pt] (1.6,1.6) rectangle (5.2,3.0);
        \node[font=\bfseries\color{charcoal}] at (3.4,2.3) {เซลล์ใส (ผลักสีประจุลบออกไป)};
        \foreach \x in {2.1, 3.4, 4.7} {
            \fill[crimsonAlert] (\x,3.6) circle (4.5pt);
            \node[white, font=\scriptsize\bfseries] at (\x,3.6) {$-$};
            \draw[->, line width=1.1pt, color=physAmber] (\x,3.2) -- (\x,3.5);
        }
        \node[font=\footnotesize, align=center, color=charcoal] at (3.4,0.6) {เช่น Nigrosin, India Ink, Eosin\\ประจุ $-$ ผลักกับผิวเซลล์ $-$ $\rightarrow$ \textbf{พื้นหลังมืด ตัวเซลล์ใส (ดูแคปซูล)}};
    \end{scope}
\end{tikzpicture}%
}
\caption{กลไกแรงทางไฟฟ้าสถิตระหว่างประจุบนผิวเซลล์แบคทีเรียกับสีย้อมประจุบวกและสีย้อมประจุลบ}
\end{figure}

\vspace{4pt}
\arucomarkerbadge{6}{7}{สนามไฟฟ้าคูลอมบ์ ศักย์ซีตาผนังเซลล์ และการยึดเกาะของสีย้อมแกรม}{จำลองประจุลบของหมู่ฟอสเฟตบนผนังเซลล์แบคทีเรียที่ดึงดูดไอออนสีย้อมประจุบวก (Crystal Violet) และผลักสีย้อมประจุลบ (Nigrosin)}{ชนิดสีย้อม (ประจุบวก / ประจุลบ), ค่า $\text{pH}$ ของตัวกลาง, ระยะห่างระหว่างประจุ $r$ ($\text{nm}$)}


\begin{worked}[การเปรียบเทียบขนาดแรงทางไฟฟ้าตามระยะห่างระหว่างสีย้อมกับผนังเซลล์]
\textbf{โจทย์ปัญหา}\quad โมเลกุลสีย้อมเมทิลีนบลูซึ่งมีประจุบวก ($+q_1$) เคลื่อนที่เข้าไปใกล้หมู่ฟอสเฟตที่มีประจุลบ ($-q_2$) บนผนังเซลล์แบคทีเรีย หากระยะห่างระหว่างประจุทั้งสองลดลงจาก $r_1 = 6\text{ nm}$ เหลือ $r_2 = 2\text{ nm}$ (เข้าใกล้ขึ้น $3$ เท่า) แรงดึงดูดทางไฟฟ้าจะเพิ่มขึ้นเป็นกี่เท่าของตอนแรก

\textbf{วิธีทำทีละขั้นตอนฉบับเข้าใจง่าย}
\begin{enumerate}
    \item \textbf{ขั้นที่ 1 ใช้หลักสัดส่วนผกผันกำลังสองจากกฎของคูลอมบ์ ($F_e \propto 1/r^2$)}
    \[
        \frac{F_2}{F_1} = \left(\frac{r_1}{r_2}\right)^2
    \]
    \item \textbf{ขั้นที่ 2 แทนค่าระยะห่างแล้วคำนวณกำลังสอง}
    \[
        \frac{F_2}{F_1} = \left(\frac{6\text{ nm}}{2\text{ nm}}\right)^2 = 3^2 = 9\text{ เท่า}
    \]
\end{enumerate}
\textbf{คำตอบ}\quad แรงดึงดูดทางไฟฟ้าระหว่างสีย้อมกับผนังเซลล์จะเพิ่มสูงขึ้นเป็น $9$ เท่า
\end{worked}

\section{ปฏิบัติการประจำสัปดาห์ที่ 8}

\deepskinfig{figures/sec_3d/ch07_s3.jpg}{ชุดปฏิบัติการที่ 7 สำหรับทดสอบแรงกระทำระหว่างประจุไฟฟ้าและการยึดเกาะของสีย้อมประจุบวกและประจุลบบนเส้นใยเซลลูโลส}

\begin{labbox}[ปฏิบัติการที่ 7\quad ไฟฟ้าสถิตและกลไกประจุของสีย้อมเซลล์จุลินทรีย์ (3 ชั่วโมง)]
\noindent\textbf{ผลลัพธ์การเรียนรู้ที่ประเมิน}\quad CLO1 (กฎของคูลอมบ์และประจุไฟฟ้า), CLO3 (ความปลอดภัยและวินัยในห้องแล็บ), CLO4 (รายงานกลุ่ม)

\vspace{4pt}
\noindent\textbf{1. วัสดุและอุปกรณ์ (ดังแสดงใน\mbox{ภาพที่} 7.4)}
\begin{itemize}
    \item ชุดทดลองไฟฟ้าสถิต (แท่งพลาสติก แท่งแก้ว ผ้าขนสัตว์ และลูกพิธหรือแผ่นอิเล็กโทรสโคป)
    \item กระดาษกรองเซลลูโลส (ซึ่งมีประจุผิวเป็นลบจำลองผนังเซลล์แบคทีเรีย), สารละลายสีย้อมประจุบวก (Methylene Blue หรือ Crystal Violet) และสีย้อมประจุลบ (Eosin หรือสีผสมอาหารประจุลบ)
\end{itemize}

\noindent\textbf{2. ขั้นตอนการทดลอง}
\begin{enumerate}
    \item \textbf{ตอนที่ 1 การทดสอบแรงดึงดูดและแรงผลักระหว่างประจุไฟฟ้า}\quad ถูแท่งพลาสติกด้วยผ้าขนสัตว์แล้วนำเข้าใกล้ลูกพิธ สังเกตการดึงดูดและการผลักกันเมื่อมีประจุชนิดเดียวกัน
    \item \textbf{ตอนที่ 2 การยึดเกาะทางไฟฟ้าสถิตของสีย้อมประจุบวกและประจุลบ}\quad หยดสีย้อมประจุบวก (Methylene Blue) และสีย้อมประจุลบ (Eosin) ปริมาตร $10\text{ }\mu\text{L}$ ลงบนแถบกระดาษกรองเซลลูโลส ทิ้งไว้ 2 นาที จากนั้นจุ่มปลายกระดาษลงในน้ำกลั่นเพื่อให้เกิดการชะล้าง (\textbf{Washing/Chromatography}) สังเกตว่าสีใดยึดเกาะแน่นกับเส้นใยประจุลบอยู่กับที่ และสีใดถูกน้ำชะเคลื่อนที่หลุดออกไปได้ง่าย
\end{enumerate}

\noindent\textbf{3. ตารางบันทึกผลการทดลอง}
\begin{center}
\footnotesize\setlength{\tabcolsep}{3.5pt}
\begin{tabularx}{\linewidth}{>{\raggedright\arraybackslash}X c c >{\raggedright\arraybackslash}X}
\toprule
\textbf{ชนิดสีย้อมที่ทดสอบ} & \textbf{ประจุของไอออนสี} & \textbf{ระยะที่ถูกน้ำชะเคลื่อนที่ ($\text{cm}$)} & \textbf{สรุปแรงกระทำทางไฟฟ้ากับเส้นใยประจุลบ} \\
\midrule
Methylene Blue & ประจุบวก ($+$) & ............ & \dotfill \\
Eosin / Nigrosin & ประจุลบ ($-$) & ............ & \dotfill \\
\bottomrule
\end{tabularx}
\end{center}
\end{labbox}

\section*{คำถามท้ายบทที่ 7 (ทบทวนแนวคิดและการคำนวณเบื้องต้น)}
\begin{enumerate}
    \item เพราะเหตุใดในการส่องดูแคปซูล (Capsule) ของแบคทีเรียด้วยเทคนิค Negative Staining จึงนิยมใช้สีย้อมประจุลบ เช่น Nigrosin แทนการใช้สีย้อมประจุบวกทั่วไป
    \item หากประจุสองตัววางห่างกันเพิ่มขึ้นเป็น $2$ เท่าของระยะเดิม ขนาดของแรงทางไฟฟ้าระหว่างประจุทั้งสองจะลดลงเหลือกี่เท่าของค่าเดิม
\end{enumerate}
